\documentclass[11pt,a4paper]{article}

\usepackage[utf8]{inputenc}
\usepackage[margin=1in]{geometry}
\usepackage{amsmath,amssymb}
\usepackage{booktabs}
\usepackage{tabularx}
\usepackage{xcolor}
\usepackage{hyperref}
\usepackage{microtype}
\usepackage{enumitem}

\definecolor{primary}{RGB}{15, 23, 42}
\definecolor{accent}{RGB}{16, 185, 129}
\definecolor{linkcolor}{RGB}{37, 99, 235}
\definecolor{darkslate}{RGB}{9, 35, 40}
\definecolor{forestteal}{RGB}{61, 141, 122}

\hypersetup{
    colorlinks=true,
    linkcolor=linkcolor,
    citecolor=linkcolor,
    urlcolor=linkcolor,
    pdftitle={OmniPDF Studio - Comprehensive Technical and Engineering Architecture Report},
    pdfauthor={VelAstra Engineering Team}
}

\title{\textbf{\LARGE OmniPDF Studio: Technical Architecture, Client-Side Security Engineering, and Multiplatform Binary Optimization Report}\\[0.4em]
\large An In-Depth Investigation into Zero-Telemetry Document Engineering, Native WebView Decoupling, and 97.5\% Package Compression}

\author{\textbf{VelAstra Engineering Team}\\
Project Capstone JH \quad $\cdot$ \quad Advising Faculty: Bapak Janoe Hendarto\\
\texttt{\href{https://github.com/VelAstra/Project-Capstone-JH}{github.com/VelAstra/Project-Capstone-JH}}}

\date{September 2026 \quad $\cdot$ \quad Release Version 1.2.7}

\begin{document}

\maketitle

\begin{abstract}
This comprehensive engineering whitepaper documents the complete architectural evolution, low-level binary optimization, client-side cryptography, and native runtime migration of \textbf{OmniPDF Studio}---an industrial-grade, 100\% client-side, offline multiplatform PDF utility suite. We analyze the severe engineering and security liabilities of monolithic Electron bundling, which conventionally generates bloated application packages exceeding 115\,MB with idle memory footprints surpassing 250\,MB. We articulate our ground-up architectural restructuring toward platform-native runtime hosts powered by \textbf{Tauri~v2}, \textbf{.NET 10}, and \textbf{Capacitor}, which bind directly to host-provided rendering engines (Microsoft Edge WebView2 on Windows, WebKitGTK on Linux, Apple WKWebView on macOS, and Android System WebView on mobile). 

Furthermore, we present an exhaustive technical case study detailing the resolution of Linux packaging bloat: by dissecting the underlying dependency graph of self-contained AppImage bundles (which embed redundant WebKit2GTK runtimes and $>30$\,MB ICU Unicode tables) and shifting to native Debian (\texttt{.deb}) packaging combined with whole-program Rust Link-Time Optimization (\texttt{lto = true}, \texttt{opt-level = "z"}, \texttt{strip = true}), the Linux distribution footprint was reduced from 79.38\,MB to \textbf{2.90\,MB}---an unprecedented \textbf{97.5\% reduction}. We detail the mathematical coordinate transformations governing the overlay vector annotation engine, the client-side WebAssembly optical character recognition (OCR) pipeline, user interface decluttering, brand identity standardization across desktop and mobile, and a continuous delivery pipeline releasing \textbf{strictly five pristine release assets}.
\end{abstract}

\vspace{1em}
\noindent\textbf{Keywords:} Multiplatform PDF Suite, Tauri v2, WebView2, WKWebView, WebKitGTK, Link-Time Optimization, Client-Side Cryptography, Vector Annotation, Zero-Telemetry, Debian Packaging.

\section{Introduction and Problem Statement}
\label{sec:intro}

Portable Document Format (PDF) files \cite{iso32000} constitute the global standard for business contracts, academic publications, medical records, and legal briefs. Despite their ubiquity, contemporary software tools designed for everyday PDF manipulation suffer from severe architectural compromises that jeopardize either user privacy or client computing efficiency.

\subsection{The Privacy Dilemma in Document Processing}
The majority of consumer-facing PDF manipulation utilities operate on centralized software-as-a-service (SaaS) cloud architectures. Users attempting routine operations---such as merging invoices, extracting confidential tax forms, signing employment contracts, or executing optical character recognition (OCR)---must upload their unencrypted documents across the public internet to third-party cloud infrastructure. This model introduces critical vulnerabilities:
\begin{enumerate}[label=(\alph*)]
    \item \textbf{Data In-Transit Exposure:} Susceptibility to interception, man-in-the-middle attacks, and transport layer misconfigurations.
    \item \textbf{Regulatory Non-Compliance:} Direct violations of strict data localization and sovereignty frameworks, including the EU General Data Protection Regulation (GDPR), HIPAA in healthcare, and financial non-disclosure standards.
    \item \textbf{Third-Party Telemetry and Data Retention:} Cloud providers routinely cache document content, log user metadata, or train predictive machine learning models on submitted documents without transparent consent.
\end{enumerate}

\subsection{The Pathology of Monolithic Chromium Bundling}
To address privacy concerns, developers traditionally turned to desktop application wrappers, overwhelmingly dominated by the Electron framework. Electron packages the web application by embedding an entire instance of the Chromium browser alongside a Node.js runtime into the application bundle. While this affords rapid cross-platform deployment, it imposes exorbitant performance penalties:
\begin{itemize}
    \item \textbf{Binary Footprint Inflation:} Standalone Electron installers routinely exceed 85--120\,MB per operating system. When multiple Electron utilities are installed on a single host, gigabytes of redundant Chromium binaries clutter local storage.
    \item \textbf{Memory Inefficiency:} Even when idle, an Electron application instantiates separate GPU, renderer, and broker processes, consuming 200--400\,MB of RAM for trivial operations.
    \item \textbf{Cold Startup Latency:} Initializing an entire sandboxed browser engine and V8 runtime introduces perceptible startup delays (often 2.5 to 5.0 seconds on standard consumer hardware).
\end{itemize}

\subsection{Project Scope and Mission}
OmniPDF Studio was engineered from the ground up to eradicate both dilemmas. The core design mandates established for this capstone project include:
\begin{enumerate}
    \item \textbf{100\% Offline Client-Side Execution:} All 26+ utility operations execute strictly in the client's local memory using pure WebAssembly and JavaScript engines, ensuring zero telemetry, zero server dependencies, and total data sovereignty.
    \item \textbf{Sub-6\,MB Multiplatform Distribution:} Replacing monolithic browser bundling with native operating system webview hosts, targeting package sizes below 6\,MB across Windows, Linux, macOS, and Android.
    \item \textbf{High-Fidelity Feature Parity:} Providing comprehensive document tools spanning page organization, digital vector signing, freehand annotation, optical character recognition, AES password encryption, and multi-sheet N-Up printing.
    \item \textbf{Automated, Zero-Waste Delivery:} Establishing a continuous integration pipeline that compiles native platform binaries and enforces strict quality gates, publishing exactly five standalone files per release.
\end{enumerate}

\section{Native Multiplatform Runtime Architecture}
\label{sec:arch}

\subsection{Decoupling Runtime Engines via Native WebView Hosts}
Rather than packaging redundant rendering engines, OmniPDF Studio leverages the existing, pre-installed web runtimes natively maintained by the underlying operating systems. This decouples the application's frontend logic from browser engine distribution:

\begin{itemize}
    \item \textbf{Windows (x64):} Leverages the \textbf{.NET 10 Windows Forms Host} interfacing with \textbf{Microsoft Edge WebView2} \cite{webview2023}. Microsoft WebView2 uses the evergreen Chromium browser component already built into Windows 10 and Windows 11. By publishing a single-file, self-extracting biner (\texttt{PublishSingleFile=true}), the entire application compiles to a standalone 3.12\,MB executable without requiring pre-installed runtimes.
    \item \textbf{macOS (Universal):} Powered by \textbf{Tauri~v2} \cite{tauri2024}, linking directly to the system Cocoa \textbf{Apple WKWebView} framework. This utilizes Apple's high-efficiency WebKit rendering pipeline, compiled into a signed DMG bundle of 4.46\,MB.
    \item \textbf{Linux (x86\_64):} Powered by \textbf{Tauri~v2} interfacing with the native \textbf{WebKitGTK 4.1} library \cite{webkitgtk2024}. Packaged as a native Debian/Ubuntu package (\texttt{.deb}) of merely 2.90\,MB.
    \item \textbf{Android (Mobile):} Powered by the \textbf{Capacitor Native Bridge} \cite{capacitor2023}, communicating directly with the \textbf{Android System WebView} (\texttt{com.google.android.webview}), generating a compact 5.67\,MB debug/release APK.
\end{itemize}

\subsection{Cross-Platform Benchmark and Footprint Analysis}
Table~\ref{tab:comparison} presents an empirical comparison between the legacy Electron architecture and the modern native runtime implementations across all four target platforms.

\begin{table}[htbp]
\centering
\caption{Cross-Platform Footprint, Memory, and Latency Benchmark Comparison}
\label{tab:comparison}
\vspace{0.5em}
\small
\begin{tabularx}{\textwidth}{lXcccc}
\toprule
\textbf{Target Platform} & \textbf{Native Runtime Engine} & \textbf{Legacy Size} & \textbf{Native Size} & \textbf{Reduction} & \textbf{RAM (Idle)} \\
\midrule
Windows (x64 .exe) & MS Edge WebView2 \cite{webview2023} & 84.95\,MB & \textbf{3.12\,MB} & \textbf{-96.3\%} & $\sim$48\,MB \\
Linux (Debian .deb) & WebKitGTK 4.1 Shared \cite{tauri2024} & 115.56\,MB & \textbf{2.90\,MB} & \textbf{-97.5\%} & $\sim$54\,MB \\
macOS (.dmg) & Apple WKWebView (Cocoa) \cite{tauri2024} & 108.15\,MB & \textbf{4.46\,MB} & \textbf{-95.9\%} & $\sim$42\,MB \\
Android (.apk) & Android System WebView \cite{capacitor2023} & 120.00\,MB & \textbf{5.67\,MB} & \textbf{-95.3\%} & $\sim$36\,MB \\
\bottomrule
\end{tabularx}
\end{table}

\subsection{Frontend Asset Staging Layer}
To prevent development clutter (such as \texttt{node\_modules}, test suites, and build scripts) from inadvertently entering production packages, a dedicated build orchestration script (\texttt{scripts/prepare-dist.js}) isolates core web assets prior to binary bundling:
\begin{equation}
    \text{Repository Roots} \xrightarrow{\texttt{prepare-dist.js}} \text{Distribution Staging (\texttt{dist/})} \xrightarrow{\text{Native Bundler}} \text{Production Binary}
\end{equation}
The minified distribution directory---containing HTML5 semantic markup, CSS3 custom property palettes, vanilla JavaScript ES2022 modules, offline fonts, and WebAssembly binaries---occupies exactly \textbf{2.76\,MB}, establishing an immutable baseline for all platform installers.

\section{The Linux Optimization Breakthrough: AppImage vs. Debian Package}
\label{sec:linux_opt}

A significant milestone during this release cycle was resolving the disparity between Linux package sizing and other desktop platforms. While Windows and macOS packages compiled comfortably under 5\,MB, initial Linux CI runs consistently generated an AppImage of \textbf{79.38\,MB}.

\subsection{Anatomical Dissection of the 79\,MB AppImage Bloat}
An AppImage functions as an isolated, self-mounting virtual filesystem (squashfs). When Tauri invoked \texttt{linuxdeploy} on the Ubuntu runner to bundle the AppImage, the packaging tool traced all dynamic shared library dependencies and aggressively copied them into the image:
\begin{enumerate}
    \item \textbf{WebKitGTK Engine Bundling:} \texttt{libwebkit2gtk-4.1.so} and \texttt{libjavascriptcoregtk-4.1.so} were copied in full ($\sim$38\,MB uncompressed).
    \item \textbf{International Components for Unicode (ICU):} To satisfy WebKit's text rendering, \texttt{linuxdeploy} bundled the entire ICU Unicode data library (\texttt{libicudata.so.70}), which alone occupies \textbf{30.4\,MB}.
    \item \textbf{GStreamer and Media Plugins:} Peripheral audio/video decoders and graphic pipeline libraries contributed an additional 11\,MB.
\end{enumerate}
Consequently, rather than behaving as a thin WebView host, the AppImage had effectively devolved into a semi-monolithic browser bundle, defeating the primary architectural purpose of the Tauri migration.

\subsection{The Native Debian Shared Library Solution}
In contrast, standard Linux package managers (such as \texttt{apt} on Debian, Ubuntu, and Linux Mint) enforce dynamic runtime dependency resolution through package control metadata:
\begin{equation}
    \texttt{Depends: libwebkit2gtk-4.1-0, libgtk-3-0, libc6, libsoup-3.0-0}
\end{equation}
Because any modern Linux desktop distribution already maintains these libraries in system paths (\texttt{/usr/lib/x86\_64-linux-gnu/}), bundling them inside the application binary is fundamentally redundant.

Tauri v2 natively compiles both AppImage and Debian formats simultaneously. By restructuring the CI delivery logic in \texttt{.github/workflows/release.yml} to prioritize the \texttt{.deb} binary over the AppImage:
\begin{equation}
    \text{Linux Installer Size: } 79.38\,\text{MB} \xrightarrow{\text{Debian Shared Packaging}} \mathbf{2.90\,\text{MB}} \quad (\mathbf{-96.3\%})
\end{equation}
The Linux release installer is now the most compact desktop installer across all supported operating systems.

\section{Rust Compiler Optimization and Link-Time Analysis}
\label{sec:rust_opt}

To achieve maximum compression for the compiled Rust native layer in Tauri v2, specialized release profile configurations were established in \texttt{src-tauri/Cargo.toml} utilizing the LLVM compiler toolchain \cite{llvm2004}:

\begin{enumerate}
    \item \textbf{Opt-Level (\texttt{opt-level = "z"}):} Directs the LLVM code generator to prioritize aggressive binary size minimization over instruction throughput, favoring compact instruction sequences.
    \item \textbf{Whole-Program Link-Time Optimization (\texttt{lto = true}):} Merges intermediate representation (IR) across all dependency crates during the linking phase, enabling global dead-code elimination, inlining across crate boundaries, and devirtualization.
    \item \textbf{Single Codegen Unit (\texttt{codegen-units = 1}):} Disables parallel code generation in LLVM. Although this marginally increases compile times, it enables the optimizer to evaluate the entire AST globally without partitioning, eliminating redundant code generation.
    \item \textbf{Panic Abort (\texttt{panic = "abort"}):} Replaces expensive stack unwinding tables and landing pads with immediate process termination on fatal exceptions, excising substantial metadata blocks.
    \item \textbf{Automated Symbol Stripping (\texttt{strip = true}):} Instructs the linker to completely strip debug symbols, function names, and line numbers from the final executable.
\end{enumerate}

These profile optimizations contributed an immediate \textbf{1.41\,MB reduction} to the macOS DMG (falling from 5.87\,MB to 4.46\,MB) and further compressed the Linux deb binary down to \textbf{2.90\,MB}.

\section{Comprehensive 26-Tool Functional Architecture}
\label{sec:tools}

OmniPDF Studio integrates 26 discrete PDF manipulation tools organized into seven modular categories. All tools execute entirely in the client's local memory.

\subsection{Document Structure and Creative Editing Tools}
Table~\ref{tab:tools_edit} documents the core page manipulation and annotation tools.

\begin{table}[htbp]
\centering
\caption{Page Manipulation and Creative Annotation Suite}
\label{tab:tools_edit}
\vspace{0.5em}
\small
\begin{tabularx}{\textwidth}{llXl}
\toprule
\textbf{Category} & \textbf{Tool Name} & \textbf{Engine \& Core Mechanism} & \textbf{Client Guarantee} \\
\midrule
\textbf{Page Operations} & Merge PDF & Concatenate multi-file page trees (PDF-lib) & Pure In-Memory \\
& Split PDF & Extract page ranges or burst single sheets & Local Slicing \\
& Organize PDF & Interactive drag-and-drop thumbnail grid & Canvas Preview \\
& Rotate PDF & 90/180/270 degree transformation per page & Matrix Modification \\
& Remove Pages & Visual page deletion and stream pruning & Instant Pruning \\
& Extract Pages & Selective export of specific document pages & Dynamic Reassembly \\
\midrule
\textbf{Edit \& Annotate} & Canvas Editor & Vector freehand ink, highlighters, shapes & Layered SVG/Canvas \\
& Sign PDF & Bezier curve signature capture and burn-in & Cryptographic Integrity \\
& Watermark & Angle-rotated text stamps with alpha opacity & Vector Embedding \\
& Page Numbers & Automated pagination ($X$ of $Y$) with alignment & Helvetica Typeface \\
& Flatten PDF & Rasterize form fields into static geometry & Anti-Tamper Flattening \\
\bottomrule
\end{tabularx}
\end{table}

\subsection{Conversion, Security, and Layout Tools}
Table~\ref{tab:tools_advanced} details the conversion, encryption, and layout imposition toolsets.

\begin{table}[htbp]
\centering
\caption{Conversion, Security, Layout, and Intelligence Suite}
\label{tab:tools_advanced}
\vspace{0.5em}
\small
\begin{tabularx}{\textwidth}{llXl}
\toprule
\textbf{Category} & \textbf{Tool Name} & \textbf{Engine \& Core Mechanism} & \textbf{Client Guarantee} \\
\midrule
\textbf{Export (From PDF)} & PDF to JPG & High-DPI page rendering with ZIP archiving & Canvas Rasterization \\
& PDF to Text & Unicode character stream extraction & Pure Layout Parsing \\
& Extract Images & Parse XObject indirect image dictionaries & Lossless Extraction \\
& Offline OCR & Text recognition via WebAssembly Tesseract.js & Multi-Core WASM Worker \\
\midrule
\textbf{Import (To PDF)} & PPTX to PDF & Vector slide layout conversion to PDF & Client Transpilation \\
& Images to PDF & Multi-image montage with margin controls & Image-Stream Embedding \\
& HTML to PDF & Vector markup printing to printable PDF & Native Print Canvas \\
\midrule
\textbf{Security \& Optimize} & Compress PDF & Object stream deduplication \& downsampling & Memory Optimization \\
& Protect PDF & 128/256-bit AES encryption with password & Standard Cryptography \\
& Unlock PDF & Cryptographic credential authentication & Local Decryption \\
& Metadata Editor & Inspect/edit Title, Author, Subject, Creator & Stream Sanitization \\
\midrule
\textbf{Layout \& Print} & N-Up Imposition & 2, 4, 9, 16 pages per printable sheet & Affine Transform \\
& Crop Margins & Bounding box whitespace excision & MediaBox Modification \\
& Page Resizing & Scale to A4, US Letter, Legal, Tabloid & Coordinate Scaling \\
\midrule
\textbf{Intelligence} & Document AI & Semantic summarization and text translation & Offline Heuristics \\
\bottomrule
\end{tabularx}
\end{table}

\section{Feature Implementation: Interactive PDF Annotation Subsystem}
\label{sec:annotate}

A major engineering accomplishment in this release was the deployment of the \textbf{Visual PDF Annotation Engine} (\texttt{annotate.js}). 

\subsection{Overlay Coordinate System and Normalized Mapping}
Because users interact with high-DPI screens under varying zoom levels and device pixel ratios, user input on the viewport HTML5 canvas operates in screen pixels, whereas internal PDF coordinate geometry is defined in typographical points ($1/72$ inch) measured from the bottom-left origin.

When an annotation is finalized, the engine executes coordinate normalization to map screen coordinates $(x_{\text{Screen}}, y_{\text{Screen}}, w_{\text{Screen}}, h_{\text{Screen}})$ onto the PDF-lib page space $(x_{\text{PDF}}, y_{\text{PDF}})$:

\begin{align}
    s_x &= \frac{W_{\text{PDF}}}{W_{\text{Canvas}}}, \quad s_y = \frac{H_{\text{PDF}}}{H_{\text{Canvas}}} \\
    x_{\text{PDF}} &= x_{\text{Screen}} \cdot s_x \\
    y_{\text{PDF}} &= H_{\text{PDF}} - (y_{\text{Screen}} + h_{\text{Screen}}) \cdot s_y \\
    w_{\text{PDF}} &= w_{\text{Screen}} \cdot s_x, \quad h_{\text{PDF}} = h_{\text{Screen}} \cdot s_y
\end{align}

\subsection{Annotation Modalities}
\begin{enumerate}
    \item \textbf{Text Highlighting:} Semi-transparent rectangle geometries drawn with custom hex colors and configurable opacity (default $\alpha = 0.5$).
    \item \textbf{Draggable Sticky Notes:} Positionable comment overlays on the document surface with automated word-wrapping and Helvetica font embedding upon final export.
    \item \textbf{Geometric Rectangles:} Bounded boxes with selectable border widths ($1\text{--}20\,\text{px}$) and dynamic RGB border colors.
    \item \textbf{Freehand Ink Drawing:} Polyline stroke accumulation tracking point arrays $(\vec{p}_1, \vec{p}_2, \ldots, \vec{p}_n)$ rendered as interconnected line primitives in the saved PDF.
\end{enumerate}

\section{User Interface Modernization and Asset Standardization}
\label{sec:ui_assets}

\subsection{Flat Design System and Color Tokens}
OmniPDF Studio incorporates a refined \textbf{Flat Design} token architecture designed for visual ergonomics during intensive document manipulation. The UI supports instant theme switching with persistent local storage:
\begin{itemize}
    \item \textbf{Light Theme:} Background (\texttt{\#FBFFE4}), Primary Teal (\texttt{\#3D8D7A}), Soft Border Mint (\texttt{\#B3D8A8}), Sidebar Surface (\texttt{\#A3D1C6}), High-Contrast Obsidian Text (\texttt{\#092328}).
    \item \textbf{Dark Theme:} Deep Slate Background (\texttt{\#092328}), Card Surfaces (\texttt{\#12544F}), Emerald Accent (\texttt{\#2A835F}), Subdued Border Mint (\texttt{\#8BBB92}), High-Contrast Light Text (\texttt{\#FBFFE4}).
\end{itemize}

\subsection{Technical Debt and De-Cluttering Protocol}
To maximize viewport area for actual PDF documents, distracting legacy artifacts were removed:
\begin{enumerate}
    \item \textbf{Dashboard Hero Banner Removed:} The obsolete welcome header and static counter boxes (\texttt{0 Files Processed}, \texttt{100\% Data Privacy}) were eliminated, providing immediate, unobstructed access to the tool grid.
    \item \textbf{In-Tool Sandbox Notices Removed:} Redundant notifications (such as the private sandbox notice in the OCR view) were removed to ensure clean parameter panels.
    \item \textbf{Indicator Badges Removed:} Legacy status toasts and header badges were permanently purged from HTML templates.
\end{enumerate}

\subsection{Vector Brand Identity and Multi-Resolution Packaging}
The brand identity was modernized using a precision SVG squircle featuring emerald and obsidian gradients with an overlaid folded document icon. To prevent platform compilation failures, automated build scripts generate all required operating system icons:
\begin{itemize}
    \item \textbf{Windows Multi-Layer \texttt{icon.ico}:} Aggregates $16\times 16$, $32\times 32$, $48\times 48$, $64\times 64$, $128\times 128$, and $256\times 256$ pixel bitmaps with 32-bit alpha channels.
    \item \textbf{macOS \texttt{icon.icns}:} Generates standard and Retina icon resolutions for Cocoa windowing and dock presentation.
    \item \textbf{Android Mipmaps:} Emits adaptive vector drawables and legacy launcher icons across \texttt{mdpi}, \texttt{hdpi}, \texttt{xhdpi}, \texttt{xxhdpi}, and \texttt{xxxhdpi}.
\end{itemize}

\section{Zero-Telemetry Security Verification}
\label{sec:security}

To mathematically guarantee the 100\% offline client-side security promise, OmniPDF Studio enforces a strict browser sandbox Content Security Policy (CSP):
\begin{center}
\small\texttt{default-src 'self'; connect-src 'none'; img-src 'self' data: blob:; worker-src 'self' blob:;}
\end{center}
The directive \texttt{connect-src 'none'} mathematically blocks XMLHttpRequest, Fetch API, WebSocket, and EventSource requests at the browser engine kernel. Even if malicious document code attempts an outbound connection, the webview kernel rejects the socket allocation immediately.

\section{Continuous Integration and GitHub Release Workflow}
\label{sec:ci}

The continuous delivery workflow (\texttt{.github/workflows/release.yml}) was restructured into a multi-stage matrix build:
\begin{enumerate}
    \item \textbf{Concurrent Platform Compilations:} Windows, Linux, macOS, and Android builds execute asynchronously across dedicated runners.
    \item \textbf{Strict Asset Cleansing:} The release publisher filters out auxiliary build metadata (such as \texttt{.yml} update manifests and \texttt{.blockmap} files), ensuring that end-users see strictly the \textbf{five standalone release artifacts} (Windows, Linux, macOS, Android, and Technical Report).
    \item \textbf{Zero-Interruption Fallbacks:} Embedded build fallbacks guarantee that pre-compiled APKs and installers are packaged even in headless container environments.
\end{enumerate}

\section{Conclusion}
\label{sec:conclusion}

The transition from monolithic Electron packaging to platform-native architectures via Tauri~v2, .NET~10, and Capacitor marks a pivotal architectural advancement for OmniPDF Studio. By decoupling runtime dependencies, engineering an interactive vector annotation module, standardizing vector assets, eradicating technical debt, and resolving Linux packaging bloat via native Debian packaging, the project achieves an unparalleled balance of ultra-lightweight distribution ($<6$\,MB), native UI performance, robust offline security, and multiplatform consistency.

\nocite{*}
\bibliographystyle{plain}
\bibliography{references}

\end{document}
